\Needspace{20\baselineskip}
\subsection{Recursive renaming of focussed states}
\label{sec:recursive-state-search}
The preceding quotient fixes box-body labels literally. We now extend search
to the recursive graph isomorphisms of Section~23, allowing independent changes
of carrier names inside every nested box. All constructions stay at one fixed
finite depth bound. Compatibility across depth bounds and with the paper's
set-based representation remains a separate obligation.

\begin{lemma}[Primitive labels and focus freedom]
Recursive label equivalence preserves each specified primitive symbol exactly.
It also preserves and reflects absence of focuses throughout a label, and
recursive graph isomorphisms preserve and reflect global focus freedom.
\end{lemma}
\begin{proof}
At depth zero label equivalence is equality. At positive depth a primitive can
only relate to the same primitive: the primitive--box alternatives are false.
This proves the first assertion, including the variable and operation symbols
needed by search.

For focus freedom, induct on depth. Primitive cases again reduce to equality;
mixed cases are impossible. For related boxes, take the supplied recursive
isomorphism of their bodies. In the forward direction, choose the inverse
image of each target edge under its edge bijection, apply the induction
hypothesis to the related labels, and cancel the bijection with its inverse.
In the reverse direction, apply the induction hypothesis to the image of each
source edge. Applying this label result edgewise, with the same inverse-image
argument, gives the graph-level equivalence.
\end{proof}
\lean{Graph/DeepFocus}{10}{19}{Linear.labelEquiv_primitive}
\lean{Graph/DeepFocus}{22}{43}{Linear.labelEquiv_focusFree}
\lean{Graph/DeepFocus}{45}{52}{Linear.DeepIso.focusFree_iff}

\begin{definition}[Focussed recursive isomorphism]
\termdef{focussed-recursive-isomorphism}
A focussed recursive isomorphism consists of a recursive graph isomorphism,
equality of the modes, and an equation carrying the distinguished focus edge
to the distinguished focus edge. Literal-label focussed isomorphisms embed
by treating equality of labels as recursive label equivalence.
\end{definition}
Identity uses recursive identity; inversion uses the inverse graph bijections
and the reversed mode equation; composition composes the graph maps and both
focus equations. Thus these constructions retain the distinguished data,
not only the unpointed graph.
\lean{Graph/DeepFocus}{57}{60}{Linear.Focussed.DeepIso}
\lean{Graph/DeepFocus}{65}{65}{Linear.Focussed.DeepIso.refl}
\lean{Graph/DeepFocus}{67}{72}{Linear.Focussed.DeepIso.symm}
\lean{Graph/DeepFocus}{74}{79}{Linear.Focussed.DeepIso.trans}
\lean{Graph/DeepFocus}{82}{82}{Linear.Focussed.DeepIso.ofIso}

The focus source and target correspond: apply the ordered incidence equations
to the distinguished edge and use its singleton source and target lists.
The port maps preserve input status, so these endpoint equations and mode
equality preserve and reflect both initiality and finality.
\lean{Graph/DeepFocus}{84}{87}{Linear.Focussed.DeepIso.source}
\lean{Graph/DeepFocus}{89}{92}{Linear.Focussed.DeepIso.target}
\lean{Graph/DeepFocus}{94}{97}{Linear.Focussed.DeepIso.initial_iff}
\lean{Graph/DeepFocus}{99}{102}{Linear.Focussed.DeepIso.final_iff}

\Needspace{18\baselineskip}
\subsection{Search and insertion under recursive renaming}
Let $a:G\cong_d H$ be a recursive graph isomorphism and $a_*$ its outer port
map. Wire labels still agree: internal labels agree by the graph-isomorphism
field, and boundary labels agree because boundary positions are fixed.
\lean{Graph/DeepSearchRenaming}{11}{16}{Linear.deepIso_vertexLabel}

Transport an operation view by the outer edge bijection and port map, retaining
its operation symbol. The primitive-label lemma proves the operation-label
certificate even when other edges carry independently renamed boxes. Ordered
incidence certificates follow from the graph equations, and arity follows
from preservation of list length. Transport a typed coordinate by keeping
its mode and applying $a_*$ to its vertex; wire-label preservation proves
the new typing certificate.
\lean{Graph/DeepSearchRenaming}{21}{30}{Linear.SearchOperation.deepRename}
\lean{Graph/DeepSearchRenaming}{86}{87}{Linear.TypedPosition.deepRename}

\begin{lemma}[Recursive transport of coordinate search]
Operation entry and return commute with recursive transport. A search step
between transported coordinates exists exactly when the original step does.
Transporting a typed coordinate forwards and then backwards returns it.
\end{lemma}
\begin{proof}
Entry examines the empty/nonempty argument-list alternatives, which mapping
preserves. Return uses optional lookup at index $i+1$, which commutes with
mapping. For search, induct on the rule: the variable case uses exact
preservation of the variable label and singleton source; entry uses its
commutation equation; return also maps the certified lookup at index $i$.
This proves forward transport. Apply it to the inverse recursive isomorphism
and cancel its port bijection to reflect steps. The same cancellation proves
the typed-coordinate inverse law, by equality of the mode and vertex and
then subtype extensionality.
\end{proof}
\lean{Graph/DeepSearchRenaming}{32}{35}{Linear.SearchOperation.deepRename_enter}
\lean{Graph/DeepSearchRenaming}{37}{40}{Linear.SearchOperation.deepRename_resume}
\lean{Graph/DeepSearchRenaming}{47}{62}{Linear.SearchStep.deepRename}
\lean{Graph/DeepSearchRenaming}{65}{78}{Linear.SearchStep.deepRename_iff}
\lean{Graph/DeepSearchRenaming}{90}{95}{Linear.TypedPosition.deepRename_symm}

\begin{theorem}[Recursive naturality of insertion]
Inserting a focus at $(m,v)$ in $G$ is recursively focus-preservingly
isomorphic to inserting it at $(m,a_*v)$ in $H$. Realization of typed
coordinates respects this isomorphism.
\end{theorem}
\begin{proof}
Extend the outer vertex and edge bijections by fixing the inserted singleton
names. Hosted vertices map to hosted images, and the fresh vertex is fixed.
Splitting commutes with this map because injectivity says $a_*w=a_*v$ exactly
when $w=v$; hence both sides take the fresh branch in the same case.

Retained edge labels inherit their recursive equivalence from $a$. The fresh
focus labels agree literally, so reflexivity supplies their recursive
equivalence. Retained source lists use commutation with splitting, retained
target lists use commutation with hosting, and the fresh focus equations use
the hosted and fresh port identities directly. The mode and fresh focus edge
are fixed. These fields give the focussed recursive isomorphism, and the
state-realization specialization follows by supplying its focus-free and
hole-free certificates.
\end{proof}
\lean{Graph/DeepInsertionRenaming}{10}{12}{Linear.FocusInsertion.deepRename_host}
\lean{Graph/DeepInsertionRenaming}{14}{15}{Linear.FocusInsertion.deepRename_fresh}
\lean{Graph/DeepInsertionRenaming}{17}{25}{Linear.FocusInsertion.deepRename_split}
\lean{Graph/DeepInsertionRenaming}{28}{67}{Linear.FocusInsertion.deepRename_iso}
\lean{Graph/DeepInsertionRenaming}{70}{77}{Linear.FocusInsertion.deepRename_focussed}
\lean{Graph/DeepInsertionRenaming}{86}{89}{Linear.TypedPosition.realize_deepRename}

\Needspace{18\baselineskip}
\subsection{Erasure under recursive renaming}
Let $a:f\cong_{f,d}g$ be a focussed recursive isomorphism. Its outer vertex
bijection carries the selected internal target to the selected internal target:
use target preservation, rewrite both selection certificates, and cancel the
internal-port constructor. Restrict this bijection to the undeleted vertices
and restrict the edge bijection to the non-focus edges, using the existing
complement-bijection construction.
\lean{Graph/DeepErasureRenaming}{13}{16}{Linear.Focussed.DeepIso.targetVertex}
\lean{Graph/DeepErasureRenaming}{19}{20}{Linear.Focussed.DeepIso.erasedVertices}
\lean{Graph/DeepErasureRenaming}{23}{24}{Linear.Focussed.DeepIso.erasedEdges}

\begin{lemma}[Retained ports and contraction]
Inclusion of retained ports commutes with these restricted maps. The joined
vertices correspond, and contraction commutes with recursive renaming.
\end{lemma}
\begin{proof}
The inclusion equation holds by cases on ports. Apply the injective inclusion
on the target side to compare joined vertices: their images are the focus
sources, already identified by $a$. For contraction, if the port is the
deleted target, both sides are corresponding joined vertices. Otherwise
injectivity ensures that its image is also not deleted. Apply inclusion and
use the contraction--inclusion identities on both sides to reduce to the
inclusion commutation equation.
\end{proof}
\lean{Graph/DeepErasureRenaming}{26}{29}{Linear.Focussed.DeepIso.erased_lift}
\lean{Graph/DeepErasureRenaming}{32}{37}{Linear.Focussed.DeepIso.erased_position}
\lean{Graph/DeepErasureRenaming}{40}{55}{Linear.Focussed.DeepIso.erased_contract}

The restricted maps form a recursive isomorphism of the erased graphs.
Internal wire labels and retained recursive edge-label relations are inherited
from $a$. For sources and targets, rewrite the original incidence equation
and commute contraction with the port map pointwise along its ordered list.
Nothing requires literal equality of boxed bodies.
\lean{Graph/DeepErasureRenaming}{58}{78}{Linear.Focussed.DeepIso.erased}

\Needspace{20\baselineskip}
\subsection{State search modulo recursive endpoint renaming}
\begin{definition}[Recursively renamed state search]
\termdef{recursive-state-search}
Write $s\longrightarrow_d t$ when a typed coordinate $q$ on the erased
underlying graph of $s$ satisfies the coordinate search rule from the extracted
coordinate of $s$, and the realization of $q$ is recursively
focus-preservingly isomorphic to $t$.
\end{definition}
Taking the actual realization as target gives a step by the identity
isomorphism. Composing the supplied endpoint isomorphism with another proves
closure under recursive target renaming.
\lean{Graph/DeepStateSearch}{12}{16}{Linear.State.DeepSearch}
\lean{Graph/DeepStateSearch}{19}{24}{Linear.State.deepSearch_realize}
\lean{Graph/DeepStateSearch}{27}{30}{Linear.State.deepSearch_target_iso}

\begin{theorem}[Recursive search invariance, determinism, and boundaries]
Two recursive-search successors of the same state are focussed recursively
isomorphic. Recursive search respects recursive isomorphisms in both arguments.
Final states have no such search step, and no such step reaches an initial
state.
\end{theorem}
\begin{proof}
For determinism, compare the two witnessing coordinate steps on the identical
erased source graph. Coordinate determinism identifies their results; subtype
extensionality identifies their typed coordinates. Compose the inverse of
one endpoint isomorphism with the other.

For source invariance, erasure naturality identifies the extracted coordinates:
the mode equation gives the first component and the joined-position equation
the second. Transport the successor coordinate and step through the erased
recursive isomorphism. Naturality of realization identifies the transported
realization with the old one, so compose its inverse with the old endpoint
isomorphism. Combine this with target closure; apply inverse isomorphisms to
obtain the converse and hence invariance in both arguments.

For a final source, its extracted coordinate is backtrack at an input, which
has no coordinate search successor. For an initial target, transport its
initiality backwards through the endpoint isomorphism and then through the
realization boundary law. The successor coordinate would be search at an
input, excluded by the coordinate no-initial-target theorem.
\end{proof}
\lean{Graph/DeepStateSearch}{33}{39}{Linear.State.deepSearch_deterministic}
\lean{Graph/DeepStateSearch}{62}{67}{Linear.State.coordinate_deepIso}
\lean{Graph/DeepStateSearch}{70}{80}{Linear.State.deepSearch_source_iso}
\lean{Graph/DeepStateSearch}{83}{90}{Linear.State.deepSearch_iso_iff}
\lean{Graph/DeepStateSearch}{42}{48}{Linear.State.final_no_deepSearch}
\lean{Graph/DeepStateSearch}{51}{59}{Linear.State.no_deepSearch_initial}

\begin{theorem}[Exact relation to literal-label state search]
Every literal-label state-search step is a recursive step. More precisely,
\[
 s\longrightarrow_d t\quad\Longleftrightarrow\quad
 \exists u,\ s\longrightarrow u\ \land\ u\cong_{f,d}t.
\]
\end{theorem}
\begin{proof}
A literal endpoint isomorphism embeds into a recursive one, giving the first
claim. For the forward implication of the equivalence, choose $u$ to be the
actual realization of the witnessing successor coordinate. This is a literal
state-search successor by identity, and the supplied recursive endpoint
isomorphism identifies it with $t$. Conversely, embed the literal step and
compose its endpoint witness with the supplied recursive isomorphism.
\end{proof}
\lean{Graph/DeepStateSearch}{93}{96}{Linear.State.search_deepSearch}
\lean{Graph/DeepStateSearch}{99}{106}{Linear.State.deepSearch_iff}
Thus the new relation changes the permitted endpoint representatives; it does
not introduce another local search rule. It is not a claim that literal-label
search between two fixed representatives is itself invariant under arbitrary
changes of box-body labels.

\Needspace{18\baselineskip}
\subsection{Recursive state classes and their search relation}
\begin{definition}[Recursive state class]
\termdef{recursive-state-class}
Quotient states by existence of a focussed recursive isomorphism. The identity,
inverse, and composition constructions prove this is an equivalence relation.
Write $[s]_d$ for the class of $s$. Descend recursive search to pairs of classes
using its proved two-argument invariance. Descend initiality and finality using
their recursive-isomorphism invariance.
\end{definition}
\lean{Graph/DeepStateQuotient}{10}{13}{Linear.deepStateSetoid}
\lean{Graph/DeepStateQuotient}{16}{16}{Linear.DeepStateClass}
\lean{Graph/DeepStateQuotient}{22}{22}{Linear.DeepStateClass.ofState}
\lean{Graph/DeepStateQuotient}{25}{27}{Linear.DeepStateClass.Search}
\lean{Graph/DeepStateQuotient}{39}{41}{Linear.DeepStateClass.Initial}
\lean{Graph/DeepStateQuotient}{44}{46}{Linear.DeepStateClass.Final}
Each descended proposition reduces to its representative proposition by
definition, giving the three representative equivalences.
\lean{Graph/DeepStateQuotient}{29}{30}{Linear.DeepStateClass.search_iff}
\lean{Graph/DeepStateQuotient}{48}{48}{Linear.DeepStateClass.initial_iff}
\lean{Graph/DeepStateQuotient}{50}{50}{Linear.DeepStateClass.final_iff}

\begin{theorem}[Deterministic search on recursive classes]
Recursive quotient search has at most one successor class. Final classes have
no search successor, and no search step reaches an initial class.
\end{theorem}
\begin{proof}
Eliminate the quotient arguments to representatives. Two steps from a common
representative have recursively isomorphic targets by state-level determinism;
quotient soundness equates their classes. Each boundary assertion reduces to
the corresponding representative-level boundary theorem.
\end{proof}
\lean{Graph/DeepStateQuotient}{33}{36}{Linear.DeepStateClass.search_deterministic}
\lean{Graph/DeepStateQuotient}{53}{56}{Linear.DeepStateClass.final_no_search}
\lean{Graph/DeepStateQuotient}{59}{62}{Linear.DeepStateClass.no_search_initial}

\Needspace{20\baselineskip}
\subsection{Projection and lifting between the two state quotients}
\begin{definition}[Forgetting literal body representatives]
\termdef{state-quotient-projection}
Define $\pi([s])=[s]_d$, mapping literal-label state classes to recursive state
classes. This is well-defined because every literal focussed isomorphism
embeds into a recursive one. On a specified representative the equation is
reflexive. The map is surjective: choose the literal class of any representative
of a recursive class.
\end{definition}
\lean{Graph/StateQuotientBridge}{10}{12}{Linear.StateClass.toDeep}
\lean{Graph/StateQuotientBridge}{14}{15}{Linear.StateClass.toDeep_ofState}
\lean{Graph/StateQuotientBridge}{18}{20}{Linear.StateClass.toDeep_surjective}

\begin{theorem}[Search projection and lifting]
A literal-class step $S\longrightarrow U$ gives
$\pi(S)\longrightarrow_d\pi(U)$. Conversely, for every specified literal
class $S$, a step $\pi(S)\longrightarrow_d T$ lifts to some step
$S\longrightarrow U$ with $\pi(U)=T$. Consequently,
\[
 \pi(S)\longrightarrow_d T\quad\Longleftrightarrow\quad
 \exists U,\ S\longrightarrow U\ \land\ \pi(U)=T.
\]
\end{theorem}
\begin{proof}
For projection, eliminate both literal classes to representatives and embed
the literal state-search step. For lifting, eliminate the source literal class
and target recursive class. The exact endpoint characterization supplies a
literal successor $u$ of the fixed source representative together with a
recursive isomorphism to the chosen target representative. Its literal class
is the desired $U$, and quotient soundness proves the projection equation.
The displayed equivalence combines projection and lifting.
\end{proof}
\lean{Graph/StateQuotientBridge}{23}{26}{Linear.StateClass.search_toDeep}
\lean{Graph/StateQuotientBridge}{30}{36}{Linear.StateClass.search_lift}
\lean{Graph/StateQuotientBridge}{39}{46}{Linear.StateClass.search_toDeep_iff}

Initiality and finality agree before and after projection, by quotient
elimination and their representative definitions. Absence of search successors
also agrees: a literal successor projects to a recursive successor, and a
recursive successor lifts to a literal successor. Taking the contrapositive
in each direction proves the normal-form equivalence.
\lean{Graph/StateQuotientBridge}{48}{50}{Linear.StateClass.initial_toDeep}
\lean{Graph/StateQuotientBridge}{52}{54}{Linear.StateClass.final_toDeep}
\lean{Graph/StateQuotientBridge}{57}{64}{Linear.StateClass.no_search_toDeep}

\Needspace{20\baselineskip}
\subsection{Exact-length executions across the quotients}
\begin{definition}[Search-only machines]
\termdef{search-only-machine}
On each state quotient, form a \texttt{Machine} whose step relation is its
search relation and whose final predicate is its boundary finality predicate.
The search boundary theorem proves the required final-normal law. Its
search-determinism theorem proves determinism of this machine.
\end{definition}
These machines deliberately contain only search. They are not the full UAM:
substitution and computation are absent. Their finite executions use the
already defined \texttt{Machine.Steps} relation, so the length counts search
steps exactly, without introducing a second execution datatype.
\lean{Graph/SearchQuotientExecutions}{10}{13}{Linear.StateClass.searchMachine}
\lean{Graph/SearchQuotientExecutions}{16}{19}{Linear.DeepStateClass.searchMachine}
\lean{Graph/SearchQuotientExecutions}{21}{24}{Linear.StateClass.searchMachine_deterministic}
\lean{Graph/SearchQuotientExecutions}{26}{29}{Linear.DeepStateClass.searchMachine_deterministic}

\begin{theorem}[Exact-length projection and lifting]
For every $n\in\mathbb N$, literal executions project with length unchanged,
and every recursive execution of length $n$ from $\pi(S)$ lifts from the
specified literal class $S$ with length unchanged. Thus
\[
 \pi(S)\longrightarrow_d^{n} T\quad\Longleftrightarrow\quad
 \exists U,\ S\longrightarrow^{n}U\ \land\ \pi(U)=T.
\]
\end{theorem}
\begin{proof}
Projection is induction on the execution: the zero-step execution remains
zero, and a step followed by a shorter run maps to its projected step followed
by the projected shorter run.

For lifting, induct on $n$, allowing arbitrary source and target classes.
At zero, the execution identifies the target with $\pi(S)$, so take $U=S$.
At $n+1$, lift the first recursive step from the specified source to a literal
successor $V$. Rewrite the recursive tail's source using $\pi(V)$, apply the
induction hypothesis from $V$, and prepend the lifted first step. The endpoint
equality and the exact length are retained. The equivalence combines both
constructions.
\end{proof}
\lean{Graph/SearchQuotientExecutions}{35}{40}{Linear.StateClass.steps_toDeep}
\lean{Graph/SearchQuotientExecutions}{43}{56}{Linear.StateClass.steps_lift}
\lean{Graph/SearchQuotientExecutions}{59}{65}{Linear.StateClass.steps_toDeep_iff}

\begin{theorem}[Exact-length search termination]
For each prescribed $n$, a successful search run of length $n$ exists from
$S$ if and only if one exists from $\pi(S)$.
\end{theorem}
\begin{proof}
Project or lift the run using the previous theorem. In the lifting direction,
rewrite finality along the endpoint projection equality and use reflection
of finality. In the projection direction, use preservation of finality.
\end{proof}
\lean{Graph/SearchQuotientExecutions}{68}{76}{Linear.StateClass.final_run_toDeep_iff}

\Needspace{18\baselineskip}
\subsection{Recursive-renaming and execution regression proofs}
The concrete two-hole box is globally focus-free because every edge in its
body is a hole rather than a focus. The established recursive isomorphism to
the box with its body's edge names exchanged transports this focus freedom.
These two boxed graphs were already proved to admit no literal-label
isomorphism, so the test exercises genuinely recursive label comparison.
\lean{Graph/DeepSearchExamples}{10}{12}{Linear.Examples.boxedTwoHoles_focusFree}
\lean{Graph/DeepSearchExamples}{16}{17}{Linear.Examples.renamedBoxedTwoHoles_focusFree}

The earlier concrete value step embeds into recursive search. Prepending that
step to the zero-step execution gives an execution of exactly one step from
its initial recursive class to its final recursive class. To test lifting from
a reconstructed source, identify its literal class with the original initial
class, use this recursive execution, and apply the exact-length termination
equivalence backwards. The result is a successful literal search run of exactly
one step from the reconstructed source. Finally, the final value's boundary
finality and the recursive no-search-from-final theorem exclude every recursive
search successor.
\lean{Graph/DeepSearchExamples}{20}{23}{Linear.Examples.value_deep_run}
\lean{Graph/DeepSearchExamples}{27}{33}{Linear.Examples.recovered_value_lift}
\lean{Graph/DeepSearchExamples}{36}{39}{Linear.Examples.final_value_deep_normal}
